% 第十二部：討論題與中華民國的對照。出處：課程倉 1151/seminar/my/導讀_W06_Waldron_細目次與大綱.md 第四節「可在討論時提出的問題」；
% 盤點_Waldron概念之中華民國司法援用.md；W06_1019_導讀問答.md「孤立而隔絕的少數」一則。全部屬導讀人的提問。
\chapter{討論題與中華民國的對照}

以下各題均為導讀人的提問，依各部的順序排列；中華民國的對照為導讀人所作，本文未言及中華民國。

\qaentry{第十二部之一}{關於攻擊的標的}

\begin{qaquestion}
\begin{enumerate}
\item 第 1 週老師的分法以宣告違憲後法律的效力為準：宣告違憲即失效者為強式，宣告不符合憲法而效力交由政治部門處理者為弱式。本文以法院得否在個案中拒絕適用為準\pg{1354--1355}。兩種標準在英國、紐西蘭兩例上結論相同，在美國與中華民國的定期失效宣告上是否相同？
\item 本文對事前的抽象審查與專設憲法法院各只寫一段\pg{1358--1359}。第肆部分 A 節批評的個案取向在抽象審查中本不存在，其餘論證對憲法法庭的適用程度如何？
\end{enumerate}
\end{qaquestion}

\qaentry{第十二部之二}{關於四項假設}

\begin{qaquestion}
\begin{enumerate}
\item 第一項假設是否把立法程序理想化？作者把法院的法條主義算作常態、把倉促或宗派性的立法劃為病態\pg{1386}，這一劃分的依據是兩機關的設立目的，可否反過來適用於法院？
\item 中華民國就同性婚姻先後有司法院釋字第 748 號解釋與 2018 年的公民投票。以這段經過逐項檢驗四項假設，中華民國落在核心情形之內嗎？
\end{enumerate}
\end{qaquestion}

\qaentry{第十二部之三}{關於程序面的理由}

\begin{qaquestion}
\begin{enumerate}
\item 本文說法院的簡單多數決欠缺公平性的論證\pg{1391--1392}。114 年憲判字第 1 號以「集體議事的多數決原則」說明憲法訴訟法第 30 條第 1 項的評決門檻；湯德宗大法官在釋字第 770 號以「法院判決不備理由，即無正當性可言」為判決正當性的來源。兩者如何回應 Waldron？
\item Ulysses 之喻\pg{1393--1394}：制憲者或立法者所拘束者為後代，與 Ulysses 拘束自己者不同，跨世代的拘束以何為正當性的基礎？Elster《Ulysses Unbound》序言引 Jens Arup Seip：在政治上，人們從不試圖約束自己，只試圖約束他人。
\end{enumerate}
\end{qaquestion}

\qaentry{第十二部之四}{關於孤立而隔絕的少數}

\begin{qaquestion}
釋字第 748 號解釋理由書提高審查標準的理由之一：
\begin{quote}
同性性傾向者因人口結構因素，為社會上孤立隔絕之少數，並因受刻板印象之影響，久為政治上之弱勢，難期經由一般民主程序扭轉其法律上劣勢地位。
\end{quote}
湯德宗大法官在釋字第 692 號不同意見書則以該案被歧視者「要非所謂社會上『孤立而隔絕之少數』（discrete and insular minority）」為由，不提高審查基準。用本文第柒部分的三點限制\pg{1403--1405}檢驗兩者：同性性傾向者是否合於作者所界定的 discrete 與 insular？司法上的多數是否同樣受刻板印象影響？解釋作成之後的公民投票與立法經過，在作者所說的長期效果上意味什麼？
\end{qaquestion}

\qaentry{第十二部之五}{關於本文的適用範圍：釋字第 380 號}

\begin{qaquestion}
本文限定其對象為法院對立法的審查\pg{1353--1354}。釋字第 380 號審查的客體為教育部所訂大學法施行細則第 22 條第 1 項後段與第 3 項，理由為逾越大學法或未經大學法授權，不在本文設定的對象之內。
\begin{enumerate}
\item 本號理由書指明，共同必修科目「應本大學自治之精神由法律明文規定，或循大學課程自主之程序由各大學自行訂定」。立法院若以法律明定國文為共同必修，大學以講學自由爭之，在四項假設成立時即為本文所稱的核心情形。此時應由法院或立法機關作終局決定？
\item 董翔飛、施文森、曾華松三位大法官的不同意見書謂，立法者修訂大學法時如認共同必修損及學術自由或大學自治「當可明文予以禁止」，立法過程紀錄「並無是類討論，亦足證立法者無意排除上項規定」。本文所駁的第三項調和論證\pg{1394}同樣從立法者未表示反對推論其同意，兩者可否並列檢驗？
\end{enumerate}
\end{qaquestion}

\qaentry{第十二部之六}{關於中文學界的接受}

\begin{qaquestion}
中文文獻中，本文見於蘇彥圖 2018 與林建志 2014 的註；蘇彥圖 2017 的註引的是 Law and Disagreement 與 Constitutionalism: A Skeptical View，未引本文。湯德宗《權力分立新論》卷二增訂五版（2026，課綱指定）是否處理本文？可請教老師。
\end{qaquestion}
